\documentclass[a4paper,12pt]{article}

% Podpora češtiny a kódování
\usepackage[utf8]{inputenc}
\usepackage[T1]{fontenc}
\usepackage[czech]{babel}

% Vzdušnější formátování odstavců (odstraní odsazení a přidá mezery mezi odstavci)
\usepackage{parskip}

% Matematické balíčky (nutné pro znaky jako \mathbb{N}, \pm a prostředí proof)
\usepackage{amsmath}
\usepackage{amsfonts}
\usepackage{amssymb}
\usepackage{amsthm}

% Nastavení okrajů stránky
\usepackage[margin=2.5cm]{geometry}

% Informace o dokumentu
\title{Algebraická teorie čísel - Příprava ke zkoušce}
\author{Poznámky z konzultací}
\date{\today}

\begin{document}

\maketitle

\subsection*{Otázka 1: Dokažte, že diofantická rovnice $x^4 - 3y^2 = 1$ má pouze triviální řešení $(\pm 1, 0)$}

\begin{proof}
Abychom vyřešili rovnici $x^4 - 3y^2 = 1$, potřebujeme najít řešení tvaru $x_n = m^2$, kde $n, m \in \mathbb{N}_0$ a $x_n$ pochází z řešení Pellovy rovnice $x^2 - 3y^2 = 1$ definované v Tvrzení 1.1.2.

Rozdělíme důkaz na dva případy podle parity $n$:

\textbf{1. Případ: $n$ je liché}
\vspace{0.5em}

Nechť $n = 2n_0 + 1$. Podle vzorců z Tvrzení 1.1.3 získáme rovnost:
$$(m - y_{n_0} - y_{n_0+1})(m + y_{n_0} + y_{n_0+1}) = 1$$

Z tohoto součinu nutně plyne, že $m = \pm 1$ a zároveň $y_{n_0} + y_{n_0+1} = 0$. To je ovšem nemožné (spor), protože pro složky řešení Pellovy rovnice vždy platí $y_{n_0} + y_{n_0+1} > 0$.

\vspace{1em}
\textbf{2. Případ: $n$ je sudé}
\vspace{0.5em}

Nechť $n = 2n_0$ je sudé. Podle Tvrzení 1.1.3 platí $2x_{n_0}^2 - 1 = x_{2n_0} = x_n = m^2$, z čehož vyplývá rovnice $2x_{n_0}^2 - m^2 = 1$. 

Redukcí této rovnice modulo 4 zjistíme, že $x_{n_0}$ musí být liché. 

Dle Tvrzení 1.1.3 je $x_k$ liché právě tehdy, když je jeho index sudý, můžeme tedy psát $n_0 = 2n_1$ a opětovným použitím tvrzení získáme $x_{n_0} = x_{2n_1} = 2x_{n_1}^2 - 1$.

\subsection*{Pomocná tvrzení z kapitoly 1.3 (Pythagorejské trojice a Pellovy rovnice)}

\textbf{Tvrzení 1.3.1 (Parametrizace pythagorejských trojic):} Pokud $x, y, z \in \mathbb{N}_0$ jsou celá čísla splňující rovnici $x^2 + y^2 = z^2$, která jsou po dvou nesoudělná, pak existují nesoudělná čísla $u, v \in \mathbb{N}_0$ taková, že:
$$z = u^2 + v^2 \quad \text{a} \quad \{x, y\} = \{u^2 - v^2, 2uv\}$$

\begin{proof}
Redukcí rovnice $x^2 + y^2 = z^2$ modulo 4 lze ukázat, že $z$ a právě jedno z čísel $x, y$ (řekněme $x$) musí být liché (čtverce lichých čísel jsou vždy $\equiv 1 \pmod 4$ a sudých $\equiv 0 \pmod 4$).

Rovnici upravíme na tvar $y^2 = z^2 - x^2 = (z - x)(z + x)$, což po vydělení čtyřmi dává:
$$\left(\frac{y}{2}\right)^2 = \frac{z - x}{2} \cdot \frac{z + x}{2}$$

Protože $z$ i $x$ jsou lichá čísla, jsou výrazy $\frac{z-x}{2}$ a $\frac{z+x}{2}$ celá čísla. Navíc jsou tato dvě čísla nesoudělná, protože jejich případný společný dělitel by musel dělit i jejich součet (což je $z$) a jejich rozdíl (což je $x$), což by byl spor s předpokladem nesoudělnosti $x$ a $z$.

Podle principu mocnin PP1 (Tvrzení 1.2.2) z toho plyne, že pokud se součin dvou nesoudělných čísel rovná čtverci, musí být oba činitelé na pravé straně samostatnými čtverci. Existují tedy nesoudělná čísla $u, v \in \mathbb{N}_0$ taková, že:
$$\frac{z + x}{2} = u^2 \quad \text{a} \quad \frac{z - x}{2} = v^2$$

Sečtením těchto dvou rovnic získáme $z = u^2 + v^2$, jejich odečtením dostaneme $x = u^2 - v^2$. Dosazením těchto vyjádření zpět do upravované rovnice pak triviálně dostáváme $\left(\frac{y}{2}\right)^2 = u^2 v^2 \implies y = 2uv$.
\end{proof}

\vspace{1em}
\textbf{Důsledek 1.3.2:} Pokud $x, y \in \mathbb{N}_0$ splňují rovnici $2x^2 - y^2 = 1$, pak existují čísla $a, b \in \mathbb{N}_0$ taková, že $a^2 - 2b^2 = 1$ a zároveň platí:
$$x \in \{2a^2 - 1 + 2ab, 2a^2 - 1 - 2ab\}$$

\begin{proof}
Z rovnice $2x^2 - y^2 = 1$ vyplývá, že $y$ musí být liché číslo. Můžeme ho tedy zapsat jako $y = 2y_0 + 1$ pro nějaké $y_0 \in \mathbb{N}_0$. Dosadíme toto vyjádření do původní rovnice:
$$2x^2 = (2y_0 + 1)^2 + 1 = 4y_0^2 + 4y_0 + 2 \implies x^2 = 2y_0^2 + 2y_0 + 1$$

Pravou stranu rovnice nyní můžeme velmi elegantně přepsat jako součet dvou po sobě jdoucích čtverců:
$$x^2 = y_0^2 + (y_0 + 1)^2$$

Dostali jsme geometrický vztah $y_0^2 + (y_0 + 1)^2 = x^2$. Protože $y_0$ a $y_0 + 1$ jsou dvě bezprostředně po sobě jdoucí celá čísla, jsou zjevně nesoudělná. Celý výraz tak tvoří primitivní pythagorejskou trojici. Podle Tvrzení 1.3.1 proto musí existovat nesoudělná čísla $u, v \in \mathbb{N}_0$ taková, že platí:
$$x = u^2 + v^2 \quad \text{a} \quad \{y_0, y_0 + 1\} = \{2uv, u^2 - v^2\}$$

Tato množinová rovnost se rozpadá na dva symetrické případy, podle toho, který prvek odpovídá kterému:

\textbf{Případ 1:} $y_0 = 2uv$ a $y_0 + 1 = u^2 - v^2$. \\
Odečtením první rovnice od druhé dostaneme vztah: $1 = (u^2 - v^2) - 2uv = (u - v)^2 - 2v^2$. \\
Zavedeme substituci $a = u - v$ a $b = v$ (z čehož vyjádřením dostaneme $v = b$ a $u = a + b$). Tím získáváme Pellovu rovnici $a^2 - 2b^2 = 1$. \\
Nyní dosadíme tato vyjádření do vztahu pro $x$: $x = u^2 + v^2 = (a + b)^2 + b^2 = a^2 + 2ab + 2b^2$. Protože z Pellovy rovnice víme, že $2b^2 = a^2 - 1$, po dosazení dostáváme finální tvar: $x = 2a^2 - 1 + 2ab$.

\textbf{Případ 2:} $y_0 = u^2 - v^2$ a $y_0 + 1 = 2uv$. \\
Odečtením první rovnice od druhé dostaneme vztah: $1 = 2uv - (u^2 - v^2) = (u + v)^2 - 2u^2$. \\
Zavedeme substituci $a = u + v$ a $b = u$ (z čehož vyjádřením dostaneme $u = b$ a $v = a - b$). Tím opět získáváme Pellovu rovnici $a^2 - 2b^2 = 1$. \\
Nyní dosadíme tato vyjádření do vztahu pro $x$: $x = u^2 + v^2 = b^2 + (a - b)^2 = a^2 - 2ab + 2b^2$. Protože opět platí $2b^2 = a^2 - 1$, po dosazení dostáváme druhý možný tvar: $x = 2a^2 - 1 - 2ab$.
\end{proof}

Důsledek 1.3.2 nám dále ukazuje, že existují čísla $a, b \in \mathbb{N}_0$ taková, že $a^2 - 2b^2 = 1$ a zároveň platí:
$$2x_{n_1}^2 - 1 = x_{n_0} = 2a^2 - 1 \pm 2ab$$

Z toho po úpravě získáme:
$$x_{n_1}^2 = a(a \pm b)$$

Protože $a$ a $b$ jsou nesoudělná, jsou nesoudělná i čísla $a$ a $a \pm b$. Podle principu mocnin PP1 (Tvrzení 1.2.2) z toho plyne, že číslo $a$ musí být čtverec. 

Podle Tvrzení 1.3.3 (řešícího rovnici $x^4 - 2y^2 = 1$) z toho vyplývá, že $a = 1$. 

Pokud $a = 1$, dostáváme $b = 0$, $x_n = x_{n_0} = 1$ a ve výsledku $x = \pm 1, y = 0$.

Tím jsme vyčerpali všechny možnosti a dokázali, že rovnice $x^4 - 3y^2 = 1$ má pouze triviální řešení $(\pm 1, 0)$.
\end{proof}

\subsection*{Otázka 2: Dokažte, že jedinými celočíselnými řešeními rovnice $x^2 - y^3 = 1$ jsou dvojice $(\pm3,2)$, $(\pm1,0)$ a $(0,-1)$}

\begin{proof}
Nechť $x, y \in \mathbb{Z}$ splňují rovnici $x^2 - y^3 = 1$. Pravou stranu rovnice vyjádříme pomocí algebraického rozkladu:
$$x^2 = y^3 + 1 = (y + 1)(y^2 - y + 1) = (y + 1)\big((y + 1)(y - 2) + 3\big)$$ 

Protože pro reálná čísla platí $y^2 - y + 1 = (y - 1/2)^2 + 3/4 \ge 0$ a z rovnice víme, že $x^2 \ge 0$, nutně dostáváme podmínku $y + 1 \ge 0$. Pro největší společný dělitel obou činitelů ze struktury výrazu vyplývá, že $\gcd(y + 1, y^2 - y + 1) \in \{1, 3\}$. Celý důkaz proto rozdělíme na dva hlavní případy:

\textbf{1. Případ: $\gcd(y + 1, y^2 - y + 1) = 1$}
\vspace{0.5em}

Pokud jsou tito dva činitelé nesoudělní, pak z jejich součinu rovného čtverci $x^2$ podle principu mocnin PP1 (Tvrzení 1.2.2) plyne, že oba činitelé musí být samostatnými čtverci. Platí tedy $y^2 - y + 1 = a^2$ pro nějaké $a \in \mathbb{N}_0$. Vynásobením celé rovnice čtyřmi získáme:
$$4y^2 - 4y + 4 = (2a)^2 \implies (2y - 1)^2 + 3 = (2a)^2 \implies 3 = (2a - 2y + 1)(2a + 2y - 1)$$ 

Číslo 3 lze nad celými čísly rozložit na součin dvou činitelů pouze jako $1 \cdot 3$ nebo $(-1) \cdot (-3)$. Vyřešením odpovídajících jednoduchých soustav rovnic pro neznámé $a$ a $y$ dostáváme pouze dvě možné dvojice: $\langle a, y \rangle = \langle \pm1, 1 \rangle$ nebo $\langle \pm1, 0 \rangle$.
\begin{itemize}
    \item Pro $y = 1$ dostáváme prvek $y + 1 = 2$, což ovšem není čtverec, což je ve sporu s předpokladem z principu PP1.
    \item Pro $y = 0$ dostáváme z původní rovnice $x^2 = 1$, což nám dává validní celočíselná řešení $(\pm1, 0)$.
\end{itemize}

\vspace{1em}
\textbf{2. Případ: $\gcd(y + 1, y^2 - y + 1) = 3$}
\vspace{0.5em}

Podle zobecněného principu mocnin PP2 (Tvrzení 1.2.3) v tomto případě existují nesoudělná čísla $a, b \in \mathbb{N}_0$ taková, že platí rovnosti:
$$y + 1 = 3a^2 \quad \text{a} \quad y^2 - y + 1 = 3b^2$$ 

Vyjádřením výrazu $(2y - 1)^2$ a úpravou druhé rovnice získáme rovnost $3(2b)^2 - (2y - 1)^2 = 3$. Z ní vidíme, že člen $2y - 1$ musí být dělitelný 3, můžeme tedy položit $2y - 1 = 3Y$ pro nějaké $Y \in \mathbb{Z}$. S zavedením substituce $X = 2b \in \mathbb{N}_0$ získáme normovaný tvar Pellovy rovnice spolu s podmínkou pro $Y$:
$$X^2 - 3Y^2 = 1 \quad \text{a} \quad Y = 2a^2 - 1$$ 

Tento systém rovnic má jedno zřejmé triviální řešení pro zápornou hodnotu $Y$, konkrétně trojici $\langle X, Y, a \rangle = \langle 2, -1, 0 \rangle$. To odpovídá hodnotě $2y - 1 = -3 \implies y = -1$, což po dosazení zpět do zadání dává řešení $(0, -1)$.

Dále již můžeme bez újmy na obecnosti předpokládat, že $Y \ge 0$. Hledáme index $n \in \mathbb{N}_0$ takový, aby platilo $Y = y_n = 2a^2 - 1$, kde $y_n$ je složka řešení Pellovy rovnice z Tvrzení 1.1.2. Protože hodnota $y_n = 2a^2 - 1$ je zjevně lichá, musí být index $n$ nutně lichým číslem. Podle explicitních vzorců z Tvrzení 1.1.3 vyjádříme $y_n$ pro liché indexy tvaru $n = 2m + 1$ jako $y_{2m+1} = 2x_m y_{m+1} - 1$. Porovnáním obou vyjádření dostáváme klíčovou rovnost:
$$2a^2 - 1 = 2x_m y_{m+1} - 1 \implies a^2 = x_m y_{m+1}$$ 

Z aritmetických vlastností řešení Pellových rovnic platí vztah $\gcd(x_m, y_{m+1}) = \gcd(x_m, x_m + 2y_m) \in \{1, 2\}$. Nyní musíme podrobně vyšetřit tyto dvě možnosti dělitelnosti:

\begin{itemize}
    \item \textbf{Pokud $\gcd(x_m, y_{m+1}) = 1$:} Podle principu mocnin PP1 (Tvrzení 1.2.2) musí být samotné číslo $x_m$ čtvercem. Z pomocné Věty 1.3.4 (kterou jsme dokázali v Otázce 1) však víme, že $x_m$ je čtvercem pouze tehdy, když $x_m = 1$, což přímo odpovídá indexu $m = 0$. Pro $m = 0$ dostáváme $n = 1$, z čehož vyplývá $Y = y_1 = 1$. Ze substitučního vztahu $2y - 1 = 3 \cdot 1$ pak získáme $y = 2$, což dosazením dává slavné netriviální řešení $(\pm3, 2)$.
    
    \item \textbf{Pokud $\gcd(x_m, y_{m+1}) = 2$:} Podle principu mocnin PP2 (Tvrzení 1.2.3) platí, že složka $y_{m+1}$ musí splňovat vztah $y_{m+1} = 2c^2$ pro nějaké $c \in \mathbb{N}_0$. Index $m+1$ musí být v tomto případě sudý (neboť $y_{m+1}$ je sudé), položme tedy $m+1 = 2k$ pro $k \in \mathbb{N}$. Podle Tvrzení 1.1.3 víme, že platí $y_{2k} = 2x_k y_k$, takže po dosazení dostáváme:
    $$2c^2 = 2x_k y_k \implies c^2 = x_k y_k$$ 
    Protože jsou čísla $x_k$ a $y_k$ prokazatelně nesoudělná, podle principu PP1 (Tvrzení 1.2.2) z toho vyplývá, že $x_k$ musí být čtverec. Podle Věty 1.3.4 to však pro jakékoliv $k \ge 1$ není možné. V tomto podpřípadě tedy žádná další řešení nevznikají.
\end{itemize}

Tím jsme vyčerpali všechny větve důkazu a ověřili, že jedinými celočíselnými řešeními rovnice $x^2 - y^3 = 1$ jsou skutečně pouze dvojice $(\pm3,2)$, $(\pm1,0)$ a $(0,-1)$.
\end{proof}

\subsection*{Otázka 3: Dokažte, že rovnice $x^m - y^2 = 1$ pro liché $m > 1$ nemá nenulové řešení}

\textbf{Věta (V. Lebesgue, 1850):} Nechť $m \ge 3$ je liché celé číslo. Diofantická rovnice $x^m - y^2 = 1$ má pouze triviální řešení $(1,0)$.

\begin{proof}
Předpokládejme pro spor, že existují celá čísla $a, b \in \mathbb{Z}$ s $b \neq 0$ taková, že $a^m - b^2 = 1$.

\textbf{1. Krok: Parita $a$ a $b$}
\vspace{0.5em}

Nejprve určíme paritu obou čísel. Kdyby bylo $b$ liché, pak by platilo $a^m = b^2 + 1 \equiv 1 + 1 \equiv 2 \pmod 4$. Jelikož je $m \ge 3$, liché číslo na $m$-tou dává modulo 4 zbytek 1 nebo 3, a sudé číslo na $m$-tou je dělitelné 4 (dává zbytek 0). Žádná taková mocnina nemůže dávat modulo 4 zbytek 2, což je spor. Proto musí být $b$ sudé (a nenulové) a číslo $a$ nutně liché.

\vspace{1em}
\textbf{2. Krok: Rozklad v oboru Gaussových celých čísel $\mathbb{Z}[i]$}
\vspace{0.5em}

Rovnici upravíme na $a^m = b^2 + 1$ a pravou stranu rozložíme v oboru Gaussových celých čísel $\mathbb{Z}[i]$:
$$a^m = (1 + bi)(1 - bi) \text{}$$

Nyní ukážeme, že činitelé $(1+bi)$ a $(1-bi)$ jsou v $\mathbb{Z}[i]$ nesoudělní. Pokud by nějaký prvek $\alpha \in \mathbb{Z}[i]$ dělil oba tyto činitele, pak by jeho norma $n = \alpha \overline{\alpha} \in \mathbb{N}$ musela v $\mathbb{Z}$ dělit součin původních činitelů, tedy liché číslo $(1+bi)(1-bi) = a^m$. Zároveň by ale $\alpha$ musela dělit jejich součet, který je roven $2$. Z toho plyne, že norma $n$ by musela dělit hodnotu $2 \cdot \overline{2} = 4$. Jediné přirozené číslo, které dělí 4 a zároveň liché číslo, je $n = 1$. Prvek $\alpha$ je tedy jednotka, což dokazuje nesoudělnost.

\vspace{1em}
\textbf{3. Krok: Aplikace principu mocnin (PP1')}
\vspace{0.5em}

Jelikož je $\mathbb{Z}[i]$ oborem s jednoznačným rozkladem (UFD), můžeme na nesoudělné činitele tvořící $m$-tou mocninu aplikovat princip mocnin PP1'. Existuje prvek $u + vi \in \mathbb{Z}[i]$ takový, že:
$$1 + bi = (u + vi)^m \quad \text{a} \quad 1 - bi = (u - vi)^m \text{}$$
K tomuto kroku jsme využili i faktu, že pro liché $m$ je každá jednotka v $\mathbb{Z}[i]$ (což jsou $\pm 1, \pm i$) sama o sobě $m$-tou mocninou, a lze ji tedy absorbovat do výrazu $(u+vi)^m$.

\vspace{1em}
\textbf{4. Krok: Zjištění reálné části $u$}
\vspace{0.5em}

Sečtením obou rovnic z předchozího kroku dostaneme:
$$2 = (u + vi)^m + (u - vi)^m = 2u \cdot \beta$$
kde $\beta \in \mathbb{Z}[i]$ pochází z rozvoje binomické věty (všechny členy obsahující $v$ bez $u$ se při sčítání vyruší). Z toho bezprostředně plyne, že $u$ musí dělit $1$, takže $u = \pm 1$.

Možnost $u = -1$ ovšem vyloučíme. Jelikož $(u^2+v^2)^m = (1+b^2)$ je liché číslo (protože $b$ je sudé), musí být $v$ sudé. Použijeme kongruenci modulo 4 v oboru $\mathbb{Z}[i]$:
$$1 + bi = (u + vi)^m \equiv u^m + m u^{m-1}vi \pmod 4 \text{}$$
Porovnáním reálných částí na obou stranách získáme $1 \equiv u^m \pmod 4$. Protože je $m$ liché, dosazení $u = -1$ by vedlo k neplatnému vztahu $1 \equiv -1 \pmod 4$. Tedy nutně $u = 1$.

\vspace{1em}
\textbf{5. Krok: Finální spor pomocí p-adického řádu}
\vspace{0.5em}

Nyní víme, že $1 + bi = (1 + vi)^m$ s reálnou částí $u = 1$. Rozvojem pravé strany a porovnáním reálných částí získáme identitu v $\mathbb{Z}$:
$$1 = \sum_{j=0}^{(m-1)/2} (-1)^j \binom{m}{2j} v^{2j} \text{}$$

První člen této sumy (pro $j=0$) je přesně $1$. Jedničky na obou stranách se odečtou a po přesunutí druhého členu s $v^2$ dostaneme:
$$-\binom{m}{2}v^2 + \sum_{j=2}^{(m-1)/2} (-1)^j \binom{m}{2j} v^{2j} = 0 \text{}$$

Z faktu, že $b \neq 0$, logicky vyplývá i to, že $v \neq 0$.
Pro $m = 3$ je výše uvedená suma zbytků prázdná, takže rovnice degeneruje na $-\binom{3}{2}v^2 = 0 \implies v = 0$, což je přímý spor s $v \neq 0$.

Pro liché $m \ge 5$ ukážeme spor přes 2-adický řád (valuaci). Označíme si první oddělený člen jako $A = \binom{m}{2}v^2$ a další členy uvnitř sumy jako $B_j = \binom{m}{2j}v^{2j}$. Upravíme $B_j$ vyjádřením přes $A$:
$$B_j = A \cdot \frac{1}{j(2j-1)}\binom{m-2}{2j-2}v^{2j-2} = A \cdot C_j \text{}$$

Jelikož je $v$ sudé číslo, platí pro 2-adický řád výrazu $C_j$ odhad:
$$ord_2(C_j) \ge 2j - 2 - \lfloor \log_2(j) \rfloor > 0 \text{}$$
Z aditivity 2-adického řádu z toho plyne, že $ord_2(B_j) = ord_2(A) + ord_2(C_j) < ord_2(B_j)$ (tedy $ord_2(A)$ je ostře menší než $ord_2(B_j)$ pro každé $j$).

Z Důsledku 2.1.2 víme, že pokud máme součet několika čísel a právě jedno z nich (zde $A$) má ostře nejmenší p-adický řád, pak takový součet nikdy nemůže být roven nule
\end{proof}

\subsection*{Otázka 4: Dokažte, že rovnice $x^2 - y^q = 1$ pro libovolné prvočíslo $q > 3$ nemá nenulové řešení}

\textbf{Věta (Chao Ko, 1965):} Nechť $q \ge 5$ je prvočíslo. Diofantická rovnice $x^2 - y^q = 1$ má pouze triviální řešení $(\pm 1, 0)$ a $(0, -1)$ [cite: 574-576].

Důkaz (podle E. Z. Cheina) provedeme tak, že si nejprve dokážeme dvě pomocná tvrzení o vlastnostech hypotetického nenulového řešení, která nás následně dovedou k přímému sporu.

\vspace{1em}
\textbf{Tvrzení 3.2.1: Pokud $q \ge 3$ je prvočíslo a $x, y \in \mathbb{Z}^*$ splňují $x^2 - y^q = 1$, pak $q$ dělí $x$.}
\vspace{0.5em}
\begin{proof}
Předpokládejme pro spor, že $q$ nedělí $x$ ($\neg(q|x)$)[cite: 542]. Rovnici si přepíšeme jako $x^2 = y^q + 1$ a pravou stranu algebraicky rozložíme:
$$x^2 = (y + 1) \cdot \frac{y^q + 1}{y + 1} \text{ [cite: 543]}$$

Z pomocného Lemmatu 3.1.1 plyne, že největší společný dělitel těchto dvou činitelů musí dělit číslo $q$[cite: 542]. Protože jsme ale předpokládali, že $q$ nedělí $x$, nemůže $q$ dělit ani $x^2$, a tudíž musí být tito dva činitelé absolutně nesoudělní[cite: 544]. 

Podle principu mocnin PP1 (Tvrzení 1.2.2) z toho plyne, že činitel $(y + 1)$ musí být čtverec, tedy $y + 1 = u^2$ pro nějaké $u \in \mathbb{N}$[cite: 545]. (Z parity řešení víme, že $y$ je sudé, a tedy $u$ je liché [cite: 545]). Z rovnice si nyní můžeme zkonstruovat dvě různá řešení Pellovy rovnice $X^2 - yY^2 = 1$:
$$x^2 - y \cdot (y^{(q-1)/2})^2 = 1 \quad \text{a} \quad u^2 - y \cdot 1^2 = 1 \text{ [cite: 546]}$$

Máme tedy řešení $\langle x, y^{(q-1)/2} \rangle$ a zjevně minimální řešení $\langle u, 1 \rangle$ [cite: 547-549]. Podle teorie Pellových rovnic (Tvrzení 1.1.1) musí existovat přirozené číslo $m$ takové, že v oboru $\mathbb{Z}[\sqrt{y}]$ platí:
$$x + y^{(q-1)/2}\sqrt{y} = (u + \sqrt{y})^m \text{ [cite: 550-551]}$$

Pokud tuto rovnost rozvineme a redukujeme modulo $y$, zjistíme, že $y$ musí dělit výraz $m u^{m-1}$ [cite: 553-554]. Protože je $y$ sudé a $u$ liché, musí být index $m$ sudé číslo[cite: 554]. 
Když pravou stranu přepíšeme pro sudé $m$ do tvaru $(u^2 + y + 2u\sqrt{y})^{m/2}$ a zredukujeme modulo $u$, zjistíme, že $u$ musí dělit $y^{(q-1)/2}$ [cite: 556-557]. 
Avšak z předpokladu $y + 1 = u^2$ víme, že $u$ a $y$ jsou nesoudělná[cite: 557]. Z toho nutně plyne $u = \pm 1$, což po dosazení dává $y = 0$ [cite: 557-558]. To je ale spor s tím, že hledáme nenulové řešení[cite: 558]. Tvrzení je tak dokázáno a $q$ musí dělit $x$[cite: 541].
\end{proof}

\vspace{1em}
\textbf{Tvrzení 3.2.2: Pokud $q \ge 3$ je prvočíslo a $x, y \in \mathbb{Z}^*$ splňují $x^2 - y^q = 1$, pak $x \equiv \pm 3 \pmod q$.}
\vspace{0.5em}
\begin{proof}
Bez újmy na obecnosti můžeme vyjádřit činitele $(x-1)$ a $(x+1)$ podle Lemmatu 3.1.2 pomocí nesoudělných čísel $a, b \in \mathbb{Z}$ ve tvaru:
$$x - 1 = 2^{q-1}a^q \quad \text{a} \quad x + 1 = 2b^q \text{ [cite: 562-563]}$$

Pomocí algebraických úprav těchto vztahů můžeme zkonstruovat rovnost:
$$b^{2q} - (2a)^q = \left(\frac{x+1}{2}\right)^2 - 2(x-1) = \left(\frac{x-3}{2}\right)^2 \text{ [cite: 563-564]}$$

Levou stranu opět rozložíme na součin:
$$(b^2 - 2a) \cdot \left(\frac{b^{2q} - (2a)^q}{b^2 - 2a}\right) = \left(\frac{x-3}{2}\right)^2 \text{ [cite: 564]}$$

Jelikož jsou čísla $2a$ a $b^2$ nesoudělná, podle Lemmatu 3.1.1 platí, že největší společný dělitel těchto dvou činitelů na levé straně dělí $q$[cite: 565]. Nastávají tak jen dvě možnosti pro tento největší společný dělitel:
\begin{itemize}
    \item \textbf{Dělitel je $q$:} Pak je výraz $(x-3)^2$ dělitelný $q$, z čehož rovnou plyne $x \equiv 3 \pmod q$ (respektive $x \equiv \pm 3 \pmod q$ před případnou změnou znaménka v Lemmatu)[cite: 566].
    \item \textbf{Dělitel je 1 (činitelé jsou nesoudělní):} Podle principu mocnin PP1 musí existovat $c \in \mathbb{N}$ takové, že $b^2 - 2a = c^2$[cite: 568]. Ze vzdálenosti nejbližších čtverců lze odvodit nerovnost $|a| \ge |b|$ [cite: 569-570]. Z druhé strany, porovnáním velikostí původních předpisů pro $a^q$ a $b^q$, platí hraniční nerovnost $|a|^q < |b|^q$, z čehož vyplývá $|a| < |b|$ [cite: 571-572]. To je přímý spor, takže tento případ nemůže nastat[cite: 572].
\end{itemize}
Tím je dokázáno, že nutně platí $x \equiv \pm 3 \pmod q$[cite: 561].
\end{proof}

\vspace{1em}
\textbf{Závěrečný krok důkazu věty Chao Ko:}
\vspace{0.5em}

Předpokládejme, že pro nějaké prvočíslo $q \ge 5$ existuje nenulové řešení $x, y \in \mathbb{Z}^*$ rovnice $x^2 - y^q = 1$[cite: 577]. 
Podle Tvrzení 3.2.1 musí platit, že $q$ dělí $x$, což v modulární aritmetice zapíšeme jako:
$$x \equiv 0 \pmod q \text{ [cite: 577]}$$

Zároveň však podle Tvrzení 3.2.2 pro stejné číslo $x$ musí platit:
$$x \equiv \pm 3 \pmod q \text{ [cite: 577-578]}$$

Když obě kongruence spojíme, dostaneme:
$$0 \equiv \pm 3 \pmod q \text{ [cite: 578]}$$

Tato podmínka říká, že prvočíslo $q$ musí dělit číslo 3[cite: 578]. Jediné prvočíslo, které dělí 3, je však samotná 3. Náš předpoklad ale od začátku striktně požaduje, aby $q \ge 5$[cite: 578]. 

Dostali jsme jasný matematický spor. Diophantická rovnice $x^2 - y^q = 1$ tedy pro žádné prvočíslo $q > 3$ nenulové celočíselné řešení nemá[cite: 578].

\end{document}